%%%%%%%%%%%%%%%%%%%%%%%%%%%%%%%%%%%%%%%%%%%%%%%%%%%%%%%%%%%%%%%%%%%%%%%%%%%%%%%%
% Template for USENIX papers.
%%%%%%%%%%%%%%%%%%%%%%%%%%%%%%%%%%%%%%%%%%%%%%%%%%%%%%%%%%%%%%%%%%%%%%%%%%%%%%%%

\documentclass[letterpaper,twocolumn,10pt]{article}
\usepackage{usenix2019_v3}
\usepackage{amsmath}
\usepackage{filecontents}

%-------------------------------------------------------------------------------
\begin{filecontents}{\jobname.bib}
%-------------------------------------------------------------------------------
@misc{bittorrentSpec,
  key = {BitTorrent Specification},
  title = {{BitTorrent} Specification},
  note = {\url{https://wiki.theory.org/BitTorrentSpecification}}
}

@misc{trackerProtocol,
  key = {BitTorrent Tracker Protocol},
  title = {{BitTorrent} Tracker Protocol},
  note = {\url{https://wiki.theory.org/BitTorrent_Tracker_Protocol}}
}

@misc{opentracker,
  key = {opentracker},
  title = {opentracker},
  note = {\url{https://erdgeist.org/gitweb/opentracker/}}
}

@misc{tnt,
  key = {tnt},
  title = {tnt: A Tiny Torrent Client},
  note = {\url{https://github.com/sujanan/tnt/}}
}

@misc{transmission,
  key = {Transmission},
  title = {Transmission BitTorrent Client},
  note = {\url{https://transmissionbt.com/}}
}

@misc{mktorrent,
  key = {mktorrent},
  title = {mktorrent},
  note = {\url{https://github.com/pobrn/mktorrent/wiki}}
}

@inproceedings{pbft,
  author = {Miguel Castro and Barbara Liskov},
  title = {Practical Byzantine Fault Tolerance},
  booktitle = {Proc. USENIX OSDI},
  month = feb,
  year = {1999},
  note = {\url{https://css.csail.mit.edu/6.824/2014/papers/castro-practicalbft.pdf}}
}

@inproceedings{hotstuff,
  author = {Maofan Yin and Dahlia Malkhi and Michael K. Reiter and Guy Golan Gueta and Ittai Abraham},
  title = {{HotStuff}: {BFT} Consensus with Linearity and Responsiveness},
  booktitle = {Proc. ACM PODC},
  month = jul,
  year = {2019},
  note = {\url{https://dl.acm.org/doi/10.1145/3293611.3331591}}
}
\end{filecontents}

%-------------------------------------------------------------------------------
\begin{document}
%-------------------------------------------------------------------------------

\date{}

\title{\Large \bf Traxos: Distributed Trackers with Paxos for Fun and Profit}

\author{
{\rm Ethan Carter Edwards}\\
Harvard University}

\maketitle

%-------------------------------------------------------------------------------
\begin{abstract}
%-------------------------------------------------------------------------------
Traxos is a BitTorrent tracker project built around a simple observation:
a tracker is a small replicated state machine. A tracker receives
``announce'' requests from peers, records which peers participate in which
torrents, and returns a peer list. This paper describes Traxos as it
currently stands: a BitTorrent tracker state machine, a simple HTTP tracker
front end, and an in-process Multi-Paxos prototype that replicates real
tracker announce commands under simulated message loss and replica failure.
The live HTTP server is not yet connected to the Paxos path, but the core
operation that must be replicated is implemented and tested as a Paxos
state-machine command.
\end{abstract}

%-------------------------------------------------------------------------------
\section{Introduction}
%-------------------------------------------------------------------------------

BitTorrent clients do not discover peers magically. In the common tracker
mode, each client sends an HTTP announce request to a tracker containing the
torrent's \texttt{info\_hash}, the client's \texttt{peer\_id}, its listening
port, transfer counters, and whether the client has started, completed, or
stopped the torrent~\cite{trackerProtocol}. The tracker replies with a list
of other peers for the same torrent. This makes the tracker operationally
important: if it is unreachable, new clients have a harder time joining the
swarm even if existing peers still have the file.

Traxos explores whether a tracker can be made highly available using the
same replicated-state-machine techniques used for databases and lock
services. The design goal is that a BitTorrent client should eventually be
able to announce to any tracker replica, and the replicas should converge
on the same peer table. If one tracker goes down, clients should be able to
continue using the others without losing recently committed tracker state.

The current system has two connected artifacts. First, I built a functional
HTTP announce server and tested it with \texttt{curl}, Transmission, and
\texttt{tnt}. Second, I extracted the tracker update logic into a
deterministic state machine and used it as the value type for a Multi-Paxos
prototype. The Paxos prototype runs in a simulator rather than as a
networked tracker cluster, but it orders and applies the same parsed
announce commands that the HTTP server uses.

%-------------------------------------------------------------------------------
\section{Background and Related Work}
%-------------------------------------------------------------------------------

The BitTorrent metainfo file contains an announce URL and an \texttt{info}
dictionary whose hash identifies the torrent~\cite{bittorrentSpec}. Tools
such as \texttt{mktorrent} create these files and set the tracker's announce
URL~\cite{mktorrent}. Clients such as Transmission and \texttt{tnt} then use
the announce URL to register themselves and request peers~\cite{transmission,tnt}.

Production trackers such as opentracker emphasize scale and low overhead for
very large public swarms~\cite{opentracker}. Traxos, however, focuses on small
deployments where availability matters more than absolute tracker throughput.
The prototype intentionally chooses a simple response format and state
representation so the protocol is easy to implement and reason about during
replication development.

Traxos uses a crash-fault-tolerant Paxos-style model rather than Byzantine
fault tolerance. Systems such as Practical Byzantine Fault Tolerance and
HotStuff tolerate arbitrary faulty replicas at the cost of stronger
cryptographic and quorum assumptions~\cite{pbft,hotstuff}. That model is
important for adversarial settings, but Traxos targets the simpler case of
tracker replicas that may crash or lose messages but are not malicious.

The tracker response format also influenced the implementation. Public
trackers commonly return a bencoded dictionary containing an
\texttt{interval} and a \texttt{peers} value. The \texttt{peers} value can
be a compact byte string or a list of dictionaries. Traxos currently returns
the list-of-dictionaries form, such as:

\begin{verbatim}
d8:intervali100e5:peersl
  d2:ip11:10.0.20.2274:porti51413ee
ee
\end{verbatim}

\noindent
Whitespace is shown here for readability; the actual response is one
continuous bencoded string. More advanced trackers implemented "compact"
response that is more efficient, but harder to implement.

%-------------------------------------------------------------------------------
\section{Design}
%-------------------------------------------------------------------------------

The central design decision is to treat tracker announces as state-machine
commands. In c++, the state type is:

\begin{verbatim}
map<info_hash, map<peer_id, peer>>
\end{verbatim}

\noindent
Each peer record stores the peer id, IPv4 address, port, uploaded and
downloaded counters, remaining bytes, and whether the peer is complete.
When an announce arrives, the tracker indexes first by \texttt{info\_hash} (the torrent)
and then by \texttt{peer\_id}. A \texttt{stopped} announce removes a peer;
other announces replace that peer's record. A peer is considered complete
when it sends \texttt{event=completed} or reports \texttt{left=0}.

This command boundary is also the Paxos boundary. Replicas do not need to agree
on parsing details; they need to agree on an ordered sequence of parsed
announce commands. In fact, some trackers and clients use UDP instead of
HTTP/TCP. Also, Traxos lacks support for the \texttt{/scrape} endpoint. In the
prototype, the Paxos request contains a serial number and a
\texttt{bt\_tracker\_announce\_request}. Applying a committed request mutates
the tracker table and produces the same tracker success response that the
single-node HTTP server would have produced.

Traxos stores \texttt{info\_hash} and \texttt{peer\_id} as 20-byte values. In
the HTTP parser these start as fixed-size \texttt{char[20]} fields, since the
BitTorrent protocol defines them as 20-byte strings. For use as map keys, the
implementation converts them to \texttt{std::array<char, 20>}; raw C arrays do
not compare as values in the way \texttt{std::map} needs.

IPv4 addresses are stored as \texttt{uint32\_t}. This avoids the awkwardness
of fixed-size C strings for addresses of different textual lengths. The
tracker parses dotted IPv4 strings into a 32-bit value using shifts, and
formats them back to dotted decimal when constructing responses.

Traxos accepts an explicit \texttt{ip} query parameter, but if the client does
not provide one it uses the source address of the TCP connection. This was
necessary for local testing: Transmission sends an \texttt{ipv4} parameter
containing the clients public IP address, which is not the address other machines
on my LAN should use (where testing was being done). Ignoring that alias and
using the TCP source address made the tracker return useful local peer IP
addresses. Similarly, \texttt{tnt} does not supply the ip parameter so it must
be inferred from the TCP source address.

The Multi-Paxos prototype uses a fixed set of replicas in an in-process
network simulator. One replica acts as leader and batches client announce
commands into proposed log suffixes. Followers promise rounds, accept
suffixes, apply commits in order, and redirect client requests to the
leader they currently believe in. If the leader becomes unreachable,
followers eventually time out, run a new prepare phase, and continue with a
higher round. The test harness deliberately allows a small amount of
shutdown lag between replicas, but it rejects divergent applied prefixes.

%-------------------------------------------------------------------------------
\section{Implementation}
%-------------------------------------------------------------------------------

The tracker is implemented as a Cotamer HTTP server. The server listens on
\texttt{0.0.0.0:9000}, accepts HTTP requests, and handles two paths:
\texttt{/announce} and \texttt{/debug}. The announce path parses query
parameters, applies the announce to the peer table, and returns a bencoded
success response. The debug path returns a plain-text dump of the local
tracker state, including torrent hashes, peer ids, IP addresses, ports,
remaining byte counts, and completion status.

The code is split into three conceptual pieces. The type layer defines
BitTorrent tracker data structures. The state layer applies announces to
the peer table, converts raw hashes into comparable keys, and constructs
tracker responses. The HTTP layer parses requests and sends responses. The
Paxos test uses the same state layer directly, which prevents the
replicated path from drifting away from the tracker semantics.

I originally considered using a bencode library for all tracker responses,
but the available library in the repository was more useful as a decoder
than as an encoder. Traxos therefore builds the small response body
directly with formatted strings. This is not a general bencode encoder, but
it is adequate for the tracker response shape currently supported:

\begin{verbatim}
d8:intervali100e5:peersl...ee
\end{verbatim}

\noindent
Failure responses are currently plain text. This is intentionally temporary:
real trackers should return bencoded failure dictionaries, but plain text
made early debugging faster and is fine for a prototype.

The Paxos implementation is currently still test-local rather than a
library used by the HTTP server. It stores accepted values in memory, uses
monotonic round numbers assigned by replica id, and keeps enough accepted
suffix metadata for a new leader to recover values during prepare. The
client model is intentionally lightweight: it generates consistent tracker
lifecycles of the form \texttt{started}, update, \texttt{completed},
\texttt{stopped}. This is not a full BitTorrent swarm simulation, but it
does exercise realistic tracker state transitions.

Creating a custom client (similar to \texttt{tnt} but actually functionally
correct) is probably necessary for networked testing, rather than coercing
Transmission or \texttt{tnt} to work.

The paxos implementation is correct insofar as it is tested with various failure
models similar to pset3. Given Paxos's flexibility, not much had to be done to
maintain support when switching the client models. All failure model tests
passed the first time.

%-------------------------------------------------------------------------------
\section{Evaluation}
%-------------------------------------------------------------------------------

Initial testing used \texttt{curl} against local announce URLs. This quickly
exposed a common source of tracker bugs: \texttt{info\_hash} is binary data
inside a URL, so it must be URL-encoded correctly. Once the hash was encoded
properly, repeated curl requests showed that peers were stored and returned
as expected.

The next test used a small torrent created with \texttt{mktorrent} for
\texttt{cs61hello.jpg}. I validated the torrent with Transmission's metadata
tools and then ran tracker tests across several machines on my home
network. The debug endpoint was essential here. A representative debug
response looked like:

\begin{verbatim}
torrents: 1
info_hash 45d84b...41f63b3 peers: 2
  peer ... ip 10.0.20.227 port 51413
  peer ... ip 10.0.10.175 port 51413
\end{verbatim}

Transmission testing was mostly successful, but the command-line output was
misleading. Transmission can print ``Seeding'' even when the progress
indicator is \texttt{[0.00]}; the progress value is the more important
signal. I also learned to use a fresh Transmission configuration directory
with \texttt{-g} during tests, because old resume state can make a client
appear to have files that are not present in the selected download
directory.

The hardest debugging session involved \texttt{tnt}. The tracker returned a
valid list-of-dictionaries peer response containing a Transmission seeder:

\begin{verbatim}
d8:intervali900e5:peersl
  d2:ip11:10.0.20.2274:porti51413ee
ee
\end{verbatim}

\noindent
I compared this shape with a Debian tracker response and found that Debian
also returned dictionary-form peers, not compact peers. I also tried
matching Debian's TCP \texttt{Connection: close} behavior, but that caused
\texttt{tnt} to fail with a socket-closed error. Network reachability tests
with \texttt{nc} showed that the seeder port could be reached. Finally,
\texttt{tcpdump} on the Transmission seeder showed no incoming TCP
connection attempt from \texttt{tnt}; only UDP traffic from Transmission was
visible. This suggested that the tracker had returned a usable peer, but
\texttt{tnt} did not attempt the expected TCP peer connection in this setup.
I basically gave up on \texttt{tnt} and decided to just stick to Transmission,
especially given that \texttt{tnt} is bespoke open-source software and not
a battle tested standard.

For replication testing, I used the \texttt{netsim} simulator from the
course framework. The main test program, \texttt{traxos-paxos-test},
constructs three Paxos replicas and a small set of tracker clients. Each
client submits parsed announce commands to random replicas; followers
redirect clients to the current leader, and clients retry on timeout. The
test runs multiple random seeds and reports the number of committed
announces by event type, for example:

\begin{verbatim}
33 announce: 17 started, 8 completed, 8 stopped
\end{verbatim}

\noindent
This is very similar to pset3. The checker compares all live replicas after the
simulated run. It requires their common applied prefix to match exactly,
permits a differential of 5 to allow for discrepancies when the simulation
ends, and checks that every acknowledged client command appears in the
reference replica's applied state. I ran the test with lossy replica-to-replica
channels (-l 0.01..0.5) and with failure modes that isolate the initial leader or isolate a
follower. These tests do not prove the implementation is production-ready, but
they give useful evidence that tracker announce commands are being ordered and
applied consistently by the Multi-Paxos prototype.

%-------------------------------------------------------------------------------
\section{Limitations and Future Work}
%-------------------------------------------------------------------------------

The current prototype is not yet a complete distributed tracker. Although
the Paxos path replicates real tracker announce commands in simulation,
the live HTTP server still applies announces to local memory directly. The
system also has no durable log, no restart recovery, no peer expiration, no
\texttt{/scrape} support, no IPv6 support, no compact peer response, and
only temporary plain-text failure responses. These are time/deadline
limitations rather than fundamental design problems.

The next major step is to move the Paxos implementation out of the test
harness and place it between HTTP parsing and state mutation. In that
design, \texttt{/announce} would parse the BitTorrent request into an
announce command, submit that command to a replicated log, and return the
response generated after the command commits. Each tracker replica would
apply the same commands in the same order, so any live replica could answer
future announce requests with a consistent view of the swarm.

There are also tracker-specific implementation deficiencies. Peer records
should eventually expire if clients stop announcing (currently, they are held
in memory but not served). Failure responses should be bencoded. The tracker
should probably support compact responses for clients that request
\texttt{compact=1}. Finally, the implementation should expose better
instrumentation than the current debug endpoint before being used outside a
controlled test network.

Additionally, while there was not sufficient time for this project, a more
open deployment model would be worth exploring. A trustless protocol could
allow independent participants to join the tracker replication group,
similar in idea to blockchain systems. In a straightforward Traxos
deployment, all replicas would likely be controlled by the same group and
potentially located close together, increasing the risk of correlated
failures.

Blockchain-style protocols are probably unsuitable for this tracker use
case because of their latency and throughput costs. However, Byzantine
fault-tolerant protocols such as HotStuff and Practical Byzantine Fault
Tolerance should be explored in future iterations ~\cite{hotstuff,pbft}.

%------------------------------------------------------------------------------
\section{Conclusion}
%-------------------------------------------------------------------------------

Traxos starts from a small but useful premise: a BitTorrent tracker is a state
machine. This prototype implements that state machine as a functional HTTP
tracker via a replicated state machine as Multi-Paxos. The result is not yet a
deployed replicated tracker, but it demonstrates the main design: parse
requests at the edge, order parsed announce commands through consensus, and
apply them deterministically to a peer table.

%-------------------------------------------------------------------------------
\section*{Acknowledgments}
%-------------------------------------------------------------------------------

Thanks to the CS 2620 staff (Eddie, Howard, and Tori) for the project
framework, Cotamer support, and office hours help. The debugging process also
benefited from comparing behavior with public tracker responses and real
BitTorrent clients.

%-------------------------------------------------------------------------------
\section*{Availability}
%-------------------------------------------------------------------------------

The current artifact is the Traxos repository, including the tracker server,
test torrents, and notes describing the debugging process.

%-------------------------------------------------------------------------------
\section*{AI Disclosure}
%-------------------------------------------------------------------------------

LLM/AI coding assistants were used during this project for proofreading/grammar
of the paper as well as assistance outlining and collecting my thoughts for the
first drafts. They were also used to help me sort out LaTeX and as a search
engine to find the template used for this paper.

Additionally, they were used to help migrate my old (pset3) Paxos implementation
from PancyDB to the interface I defined for the tracker implementation.

The tracker was implemented from scratch and without the assistance of any
coding assistants.


%-------------------------------------------------------------------------------
\bibliographystyle{plain}
\bibliography{\jobname}

%%%%%%%%%%%%%%%%%%%%%%%%%%%%%%%%%%%%%%%%%%%%%%%%%%%%%%%%%%%%%%%%%%%%%%%%%%%%%%%%
\end{document}
%%%%%%%%%%%%%%%%%%%%%%%%%%%%%%%%%%%%%%%%%%%%%%%%%%%%%%%%%%%%%%%%%%%%%%%%%%%%%%%%
